\documentclass[14pt]{article}
\usepackage{amsmath}
\usepackage{graphicx}
\usepackage{hyperref}
\usepackage[latin1]{inputenc}
\title{Il decadimento radioattivo}
\author{Rocca Patrick}
\begin{document}
\date{15 maggio 2021}
\maketitle

Un atomo radioattivo decade liberando particelle e onde elettromagnetiche.
Sia $N(t)$ la massa al tempo $t$, sappiamo per via sperimentale che:
$$\Delta N(t) = N(t + \Delta t) - N(t) = -\lambda N(t)\Delta t$$
La variazione dipende dalla massa $N(t)$ dal tempo di osservazione $t$ e da un valore $\lambda$ caratteristico della sostanza radioattiva.
Facendo un po\' di calcoli:
$$\Delta N(t) = -\lambda N(t)\Delta t $$
$$\frac{\Delta N(t)}{\Delta t} = -\lambda N(t)$$
Se $\Delta t \to 0$
$$\frac{dN(t)}{dt} = -\lambda N(t)$$
Abbiamo quindi un\'equazione differenziale di primo ordine del tipo $$ay\'(t) + by(t) = 0$$
Bisonga dunque determinare una funzione che soddisfi l\'equazione. Come fare?\
Sappiamo che $N\'(t)=aN(t)$, con $a = -\lambda$, possiamo dunque procedere nel seguente modo:
$$\frac{N\'(t)}{N(t)} = a$$
Ricordo che $\frac{f\'(x)}{f(x)} = D(ln[f(x)])$, quindi:
$$D(ln[N(t)]) = a$$
Se la derivata di $ln[N(t)]$ è una costante, allora $ln[N(t)]$ è una retta, dunque:
$$ln[N(t)] = at + q$$
$$N(t) = e^{at+q} = e^{at}e^{q} = ke^{-\lambda t} \text{   con } k=e^{q}$$
Rimane da determinare $k$, il che si risolve ponendo $t = 0$ e calcolando:
$$N(0) = ke^0 = k = N_0$$
Deduciamo dunque che $k$ non è altro che la massa $N$ al tempo $t = 0$, che indichiamo qunidi con $N_0$. \\
Concludiamo quindi che:
$$N(t)=N_0 e^{-\lambda t}$$
\end{document}